\chapter{Sprint 04 : Tests, Validation et Déploiement}
\minitoc
\newpage

\section{Introduction}
Ce dernier sprint est consacré à la consolidation du système. Après la phase de réalisation, il devient nécessaire de vérifier la cohérence fonctionnelle des modules, de confirmer la séparation des droits selon les rôles, de corriger les anomalies de stabilisation et de préparer une exécution locale reproductible de l'ensemble de la solution.

\section{Objectifs du Sprint}
\begin{itemize}
    \item vérifier les parcours métier essentiels sur les différents modules ;
    \item contrôler la cohérence des permissions RBAC ;
    \item stabiliser les interactions entre backend, frontend, portail client et mobile ;
    \item valider le comportement du module d'assistance documentaire ;
    \item préparer une exécution locale reproductible via Docker Compose.
\end{itemize}

\section{Tests fonctionnels}

\subsection{Dispositif de tests automatisés}
Le projet intègre plusieurs niveaux de test afin de couvrir les couches principales du système :

\begin{table}[H]
\centering
\caption{Stratégie de tests du projet}
\label{tab:strategie_tests}
\begin{tabularx}{\textwidth}{|l|l|X|}
\hline
\textbf{Niveau} & \textbf{Outil} & \textbf{Couverture visée} \\
\hline
Backend unitaire & Pytest & Authentification, logique métier et scénarios principaux du CRM \\
\hline
Backend intégration & Pytest + HTTP & Vérification de l'API, santé de l'application et parcours d'intégration \\
\hline
Frontend & Vitest & Comportements de composants, services et contextes web \\
\hline
Mobile & Jest & Vérification des services et comportements clés côté application mobile \\
\hline
Orchestration & Docker Compose + script PowerShell & Lancement coordonné de la pile d'exécution et des campagnes de tests \\
\hline
\end{tabularx}
\end{table}

En complément des tests automatisés, des scénarios fonctionnels ont été revus manuellement pour valider les flux multi-rôles les plus sensibles.

\subsection{Matrice de tests RBAC}
La matrice suivante synthétise les autorisations principales observées dans l'implémentation actuelle du système.

\begin{table}[H]
\centering
\footnotesize
\caption{Synthèse des permissions RBAC}
\label{tab:rbac}
\begin{tabularx}{\textwidth}{|X|c|c|c|c|c|}
\hline
\textbf{Opération} & \textbf{Admin} & \textbf{Manager} & \textbf{Comm.} & \textbf{Support} & \textbf{Client} \\
\hline
Connexion et accès au profil & $\checkmark$ & $\checkmark$ & $\checkmark$ & $\checkmark$ & $\checkmark$ \\
\hline
Accès au CRM interne & $\checkmark$ & $\checkmark$ & $\checkmark$ & $\checkmark$ & $\times$ \\
\hline
Création / mise à jour des leads & $\checkmark$ & $\checkmark$ & $\checkmark$ & $\checkmark$ & $\times$ \\
\hline
Conversion d'un lead qualifié & $\checkmark$ & $\checkmark$ & $\checkmark$ & $\times$ & $\times$ \\
\hline
Gestion des comptes, contacts et activités & $\checkmark$ & $\checkmark$ & $\checkmark$ & $\checkmark$ & $\times$ \\
\hline
Création d'un ticket interne & $\checkmark$ & $\checkmark$ & $\times$ & $\checkmark$ & $\times$ \\
\hline
Création d'un ticket via portail client & $\times$ & $\times$ & $\times$ & $\times$ & $\checkmark$ \\
\hline
Consultation des rapports ventes & $\checkmark$ & $\checkmark$ & $\checkmark$ & $\times$ & $\times$ \\
\hline
Consultation des rapports support & $\checkmark$ & $\checkmark$ & $\times$ & $\checkmark$ & $\times$ \\
\hline
Consultation du journal d'audit & $\checkmark$ & $\checkmark$ & $\times$ & $\times$ & $\times$ \\
\hline
Gestion des utilisateurs & $\checkmark$ & $\times$ & $\times$ & $\times$ & $\times$ \\
\hline
Utilisation du chatbot documentaire & $\checkmark$ & $\checkmark$ & $\checkmark$ & $\checkmark$ & $\checkmark$ \\
\hline
\end{tabularx}
\end{table}

\subsection{Scénario de test end-to-end}
Pour vérifier la cohérence globale du système, un scénario transverse a été retenu :
\begin{enumerate}
    \item un administrateur crée les utilisateurs internes nécessaires et rattache un utilisateur client à un compte ;
    \item un commercial crée un lead, le qualifie puis le convertit en compte, contact et deal ;
    \item un commercial crée un devis lié au deal ;
    \item le client consulte le devis depuis son portail ;
    \item le client ouvre un ticket de support ;
    \item un agent support ajoute une note interne puis une réponse publique ;
    \item le client visualise uniquement les réponses publiques et peut ajouter un nouveau message ;
    \item un manager consulte les rapports de vente et de support ;
    \item un administrateur vérifie la trace des opérations dans l'audit ;
    \item le client utilise l'assistant documentaire pour poser une question relative au support ou aux services.
\end{enumerate}

\section{Bugs identifiés et corrigés}
La phase de stabilisation a permis d'identifier plusieurs points d'amélioration concrets, touchant à la sécurité, à l'ergonomie et à la cohérence entre interface et API.

\begin{table}[H]
\centering
\footnotesize
\caption{Exemples d'anomalies de stabilisation traitées}
\label{tab:bugs}
\begin{tabularx}{\textwidth}{|c|l|X|X|c|}
\hline
\textbf{ID} & \textbf{Zone} & \textbf{Symptôme observé} & \textbf{Correction apportée} & \textbf{Statut} \\
\hline
BUG-1 & Utilisateurs & Un utilisateur authentifié pouvait consulter un autre profil sans contrôle suffisant & Restriction de l'accès au profil de l'utilisateur lui-même ou aux rôles de supervision autorisés & $\checkmark$ \\
\hline
BUG-2 & Tickets / Frontend & Les agents support ne disposaient pas d'une liste exploitable des utilisateurs assignables & Ajout d'un endpoint dédié aux utilisateurs assignables et raccordement du formulaire de ticket & $\checkmark$ \\
\hline
BUG-3 & Tickets / UI & Le formulaire de création proposait des champs non cohérents avec l'API de création & Alignement du comportement du formulaire avec la logique backend selon création ou modification & $\checkmark$ \\
\hline
BUG-4 & Tickets / Messagerie & Le nom de l'auteur n'était pas correctement restitué dans l'affichage des messages & Enrichissement de la réponse API pour inclure l'auteur lisible côté interface & $\checkmark$ \\
\hline
BUG-5 & Reporting support & Le calcul du temps moyen de résolution dépendait d'une fonction SQL spécifique au moteur MySQL & Remplacement par un calcul portable au niveau applicatif afin d'améliorer la robustesse des tests et des exécutions & $\checkmark$ \\
\hline
\end{tabularx}
\end{table}

\subsection{Méthodologie de correction}
Pour chaque anomalie identifiée, une même démarche a été appliquée :
\begin{enumerate}
    \item reproduction du problème sur l'interface ou via l'API ;
    \item isolement de la cause dans la couche concernée ;
    \item correction locale en privilégiant la cohérence avec l'architecture existante ;
    \item vérification fonctionnelle du comportement corrigé ;
    \item contrôle rapide des modules adjacents afin de limiter les régressions.
\end{enumerate}

\section{Validation du chatbot IA}
Le module d'assistance a été évalué sur des questions cohérentes avec la nature documentaire de sa base de connaissances.

\begin{table}[H]
\centering
\caption{Validation fonctionnelle du module d'assistance documentaire}
\label{tab:test_chatbot}
\begin{tabularx}{\textwidth}{|X|X|c|}
\hline
\textbf{Question posée} & \textbf{Type de réponse attendue} & \textbf{Validation} \\
\hline
Quels services propose FRS ? & Présentation synthétique des offres de l'entreprise & $\checkmark$ \\
\hline
Comment fonctionne le support client ? & Rappel des principes du guide support et du portail & $\checkmark$ \\
\hline
Quels sont les documents ou références de projets disponibles ? & Restitution des références et projets mis en avant & $\checkmark$ \\
\hline
Quelles informations trouve-t-on dans les conditions générales ? & Réponse orientée vers le document contractuel correspondant & $\checkmark$ \\
\hline
\end{tabularx}
\end{table}

\section{Résultats du Sprint}
\begin{table}[H]
\centering
\caption{Livrables du Sprint 04}
\label{tab:livrables_s4}
\begin{tabularx}{\textwidth}{|X|c|}
\hline
\textbf{Livrable} & \textbf{Statut} \\
\hline
Dispositif de tests automatisés backend, frontend et mobile structuré & $\checkmark$ \\
\hline
Matrice RBAC revue sur les parcours principaux & $\checkmark$ \\
\hline
Anomalies de stabilisation corrigées & $\checkmark$ \\
\hline
Validation du portail client et des flux tickets & $\checkmark$ \\
\hline
Validation du chatbot documentaire sur des cas d'usage réalistes & $\checkmark$ \\
\hline
Configuration de déploiement local Docker Compose finalisée & $\checkmark$ \\
\hline
\end{tabularx}
\end{table}

\section{Conclusion}
Le quatrième sprint a permis de consolider le projet et de le rendre plus robuste. La revue fonctionnelle a confirmé la cohérence des grands flux, la vérification RBAC a clarifié le périmètre de chaque rôle, et la phase de stabilisation a amélioré la sécurité ainsi que la cohérence entre interfaces et API. Le projet dispose désormais d'une base documentaire, technique et fonctionnelle nettement plus solide pour une démonstration ou une poursuite de développement.
